% 22_body.tex — 산출물 ㉒ 고객 인터뷰 스냅샷 팩 (빈 템플릿 / 완성 예시 공용 본문) — gen_product_templates.py 가 Product_specs.py 에서 생성
% 드라이버: 22_인터뷰스냅샷팩_빈템플릿.tex (\wbblanktrue) / 22_인터뷰스냅샷팩_완성예시.tex (\wbblankfalse — 워크북 §X.5 확정 후)
\documentclass[10pt,a4paper]{article}
\usepackage[a4paper,margin=16mm,top=14mm,bottom=20mm]{geometry}
\def\DocNum{22}\def\DocDay{5}\def\DocIdx{2}
\def\DocTitle{고객 인터뷰 스냅샷 팩}
\def\DocSub{Customer Interview Snapshot Pack}
\def\DocVariant{\B{빈 템플릿}{딥게이지 완성 예시}}
\usepackage{wbstyle}
\def\DocCirc{㉒}   % newunicodechar(㉛~㊵ 폴백) 뒤에 정의해야 활성 문자로 읽힌다
\begin{document}
\docheader
 {\Bg{[회사명 · 제품명]}{딥게이지 · GAUGE}}
 {\Bg{[v0.x / D+n]}{v1.0 / D+10~75 (2026-03-12~05-15)}}
 {\Bg{[인터뷰어]}{강민수(19건) · 윤하경(8건)\rd{7mm}}}
 {\Bg{[검토자 — CPO]}{윤하경 CPO\rd{7mm}}}

%% ------------------------------------------------------------------ 1
\secbar{1. 인터뷰 계획}
\guide{역할군 5(검사원·품질팀장·공장장·대표·1차사 SQE)별 목표 건수·기간·리크루팅 경로. 정본 D+90 실적은 27건(11/8/4/3/1).}
{\footnotesize
\begin{tabularx}{\linewidth}{|L{30mm}|L{20mm}|L{16mm}|Y|L{24mm}|}\hline
\hd{역할군} & \hd{목표 건수} & \hd{실제} & \hd{리크루팅 경로} & \hd{담당} \\ \hline
\ifwbblank
\rd{8mm}검사원 &  &  &  &  \\ \hline
\rd{8mm}품질팀장 &  &  &  &  \\ \hline
\rd{8mm}공장장 &  &  &  &  \\ \hline
\rd{8mm}대표 &  &  &  &  \\ \hline
\rd{8mm}1차사 SQE &  &  &  &  \\ \hline
\rd{8mm}계 &  &  &  &  \\ \hline
\else
검사원 & 10 & 11 & 파일럿 후보 공장 현장 동행 · 스마트허브 소개 & 강민수·윤하경 \\ \hline
품질팀장 & 8 & 8 & 전 직장 인맥 6 · 스마트허브 2 & 강민수 \\ \hline
공장장 & 4 & 4 & 팀장 소개 & 강민수 \\ \hline
대표 & 3 & 3 & 스마트허브 · 멘토 & 강민수 \\ \hline
1차사 SQE & 1 & 1 & 전 직장 동료 소개(H사 김도윤) & 강민수 \\ \hline
계 & 26 & 27 & 12개사 — 안산 5 · 시흥 2 · 화성 2 · 평택 1 · 충주 1 · H사 · 편향: 절삭·단조 6, 배석 4 & — \\ \hline
\fi
\end{tabularx}}

%% ------------------------------------------------------------------ 2
\secbar{2. 인터뷰 가이드}
\guide{JTBD식 질문 8~10개 — 최근에 실제로 있었던 일을 묻는다. 금지 질문("~하면 쓰시겠어요?", 기능 제시)을 따로 적는다.}
{\footnotesize
\begin{tabularx}{\linewidth}{|L{8mm}|Y|L{54mm}|}\hline
\hd{\#} & \hd{질문} & \hd{얻으려는 것} \\ \hline
\ifwbblank
\rd{8mm} &  &  \\ \hline
\rd{8mm} &  &  \\ \hline
\rd{8mm} &  &  \\ \hline
\rd{8mm} &  &  \\ \hline
\rd{8mm} &  &  \\ \hline
\rd{8mm} &  &  \\ \hline
\rd{8mm} &  &  \\ \hline
\rd{8mm} &  &  \\ \hline
\rd{8mm} &  &  \\ \hline
\rd{8mm} &  &  \\ \hline
\else
1 & 어제 하루를 처음부터 끝까지 말씀해 주세요 — 검사는 몇 시에, 기록은 언제? & 일상 속 기록의 위치와 시간 \\ \hline
2 & 가장 최근에 부적합이 나왔던 때 — 그날 무슨 일이 있었나요? & 최근 사례(이야기) \\ \hline
3 & 그때 사진은 찍으셨나요? 어디에 보냈나요? 그다음엔? & 실제 행동·도구 \\ \hline
4 & 1차사에 통보되기까지 며칠? 무엇 때문에? & G-34 검증 \\ \hline
5 & 8D는 누가 언제 썼나요? 얼마나 걸렸나요? & G-35 검증 \\ \hline
6 & 기록을 옮겨 적는 일이 있나요? 어디서 어디로, 하루에 얼마나? & G-33 검증 \\ \hline
7 & 기록을 바꿔 보려 한 적이 있나요? 어떻게 됐나요? & 시도·실패 이유 \\ \hline
8 & 감사·클레임 때 '누가 언제 무엇을 보고 판정했는지' 찾아야 했던 적은? & 판정 근거 추적 \\ \hline
9 & 1차사 제출 방식이 최근에 바뀐 적이 있나요? & 외부 압력 \\ \hline
금지 & '~하면 쓰시겠어요?' · '얼마나 불편하세요?' · 기능 제시 · 가격 언급 &  \\ \hline
\fi
\end{tabularx}}

%% ------------------------------------------------------------------ 3
\landscapeon
\secbar{3. 인터뷰 스냅샷 카드 (1건 = 1행 · 27행까지 별지 반복)}
\guide{ID·역할·회사 규모·최근 사례(이야기)·관찰된 행동·자기 보고와의 괴리·인용문·기회 후보. 관찰 열이 비면 스냅샷이 아니다.}
{\footnotesize
\begin{tabularx}{\linewidth}{|L{12mm}|L{30mm}|Y|L{44mm}|L{44mm}|L{44mm}|}\hline
\hd{ID} & \hd{역할·회사} & \hd{최근 사례(무슨 일이 있었나)} & \hd{관찰된 행동} & \hd{인용문} & \hd{기회 후보(해법 명사 없이)} \\ \hline
\ifwbblank
\rd{14mm} &  &  &  &  &  \\ \hline
\rd{14mm} &  &  &  &  &  \\ \hline
\rd{14mm} &  &  &  &  &  \\ \hline
\rd{14mm} &  &  &  &  &  \\ \hline
\rd{14mm} &  &  &  &  &  \\ \hline
\rd{14mm} &  &  &  &  &  \\ \hline
\else
I-01 & 검사원·C1 세영정공 & 초품 검사 → 종이 일지 → 퇴근 전 엑셀 전기 34분 & 일지 여백 개인 약어 · 엑셀은 팀장 PC & '적는 게 일이지 검사가 일이 아니에요' & OP-01 \\ \hline
I-02 & 검사원·C1 박선재(52, 촉탁) & 2023 태블릿 → 6개월 후 종이 복귀 & 배석 전후 발언 변화 · 오후분 종이 메모 별도 & '다 해봤어요. 결국 종이에 쓰고 나중에 옮겨 적어요' & OP-01·05 \\ \hline
I-03 & 품질팀장·C1 정미영 & 크랙 부적합 → 근거 수집 이틀 → 통보 4일째 & 카톡 사진 → PC → 엑셀 첨부 & '통보가 늦는 건 자료 모으느라' & OP-02·03 \\ \hline
I-04 & 검사원·C2 동보프레스 & 같은 부위 두 번 측정 — 기준서 두 버전 & 벽의 기준서가 2024년판 & '어느 게 맞는지 몰라서 둘 다 재요' & OP-06·01 \\ \hline
I-05 & 품질팀장·C2 & 8D 1건 13시간, 야근 이틀 & 사진 고르기 2시간 & '8D는 쓰는 게 아니라 붙이는 거예요' & OP-03 \\ \hline
I-06 & 공장장·C2 & 신입 판정 3개월 재검 & 숙련자가 뒤에서 다시 본다 & '판정이 사람 따라 다르면 시스템이 무슨 소용' & OP-05 \\ \hline
I-07 & 검사원·C3 태산정밀 & 라인 2개 겸임, 기록은 점심·퇴근 후 & 메모지에 숫자만 & '라인 돌 때는 못 적어요' & OP-01 \\ \hline
I-08 & 품질팀장·C3 & 포털 규격 분기마다 변경 & 캡처 인쇄해 벽에 & '바뀔 때마다 처음부터 배워요' & OP-07·02 \\ \hline
I-09 & 대표·C3 & 클레임 2건 4,800만, 재발 방지 근거 요구 & 서류철을 꺼내 보여 줌 & '돈보다 그때 누가 봤냐를 못 대는 게 문제' & OP-04 \\ \hline
I-10 & 검사원·C4 한빛금속 & 2023 태블릿 실패 & 태블릿이 서랍에, 배터리 없음 & '입력이 종이보다 느렸어요' & OP-01 \\ \hline
I-11 & 검사원·C4 & 오후 3시 이후 조도 & 손전등 · 수기 메모 '3시 이후' & '이 라인은 오후엔 잘 안 보여요' & OP-05·04 \\ \hline
I-12 & 품질팀장·C4 & 감사에서 판정 근거 추적 실패 & 파일 3개 폴더 분산 & '뭘 보고 그랬는지는 없어요' & OP-04 \\ \hline
I-13 & 검사원·C5 경일정공 & 검사원 1명 — 기록이 곧 검사 지연 & 검사 대기 부품 쌓임 & '제가 적는 동안 라인이 기다려요' & OP-01 \\ \hline
I-14 & 대표·C5 & MES 없음, 엑셀 12개 & 대표가 직접 취합 & '우리 규모에 시스템은 사치' & OP-01·08 \\ \hline
I-15 & 검사원·C6 우진테크(사출) & 외관 결함 — 기준 사진 없음 & 선배의 옛 사진을 폰에 & '이게 불량인지 사람마다 달라요' & OP-05·06 \\ \hline
I-16 & 품질팀장·C6 & 8D 초안 얘기에 즉답 & (회의실) & '그거 되면 제일 먼저 씁니다' & OP-03 \\ \hline
I-17 & 공장장·C6 & 데이터로 공정 개선 희망 & (회의실) & '기록부터 있어야 뭘 고치죠' & OP-08 \\ \hline
I-18 & 검사원·C7 대양정밀 — 배석 & 문제 없다고 답함 & 배석 · 발언 짧음 & '괜찮아요, 할 만해요' & OP-01(관찰 불가) \\ \hline
I-19 & 품질팀장·C7 & 통보 지연 페널티 1,200만 & 통지서 보관 & '하루만 빨랐어도' & OP-02 \\ \hline
I-20 & 검사원·C8 신광기공 & 마스터 217개 — 실제 60개 & 쓰는 행만 형광펜 & '나머지는 한 번도 안 봐요' & OP-06·01 \\ \hline
I-21 & 품질팀장·C8 & MES 있으나 엑셀 병행 & 이중 입력 & 'MES는 감사용, 엑셀은 우리용' & OP-01 \\ \hline
I-22 & 공장장·C8 & IT 담당 없음 & 서버실 없음 & '클라우드면 하고 서버는 못 둡니다' & OP-07 \\ \hline
I-23 & 대표·C10 성진금속 — 배석 & 1차사 요구 시 도입 & (회의실) & '위에서 하라면 하죠' & OP-07·04 \\ \hline
I-24 & 검사원·C9 미래오토(사출) & 사진은 이미 찍는다 — 카톡 & 검사 중 폰 촬영 → 카톡 & '사진은 찍어요. 그다음이 문제죠' & OP-01·02 \\ \hline
I-25 & 품질팀장·C9 & 카톡 → 엑셀 → 포털 세 번 & 월별 사진 폴더 & '세 번 옮기면 한 번은 틀려요' & OP-02·01 \\ \hline
I-26 & 공장장·C10 — 배석 & 자동화보다 판정 일관성 & (회의실) & '빨리보다 똑같이' & OP-05 \\ \hline
I-27 & 1차사 SQE·C11 H사 김도윤 — 회의실 & 2025 도면 유출 4,200장 → SQ 강등 & (회의실) '공문으로' & '도면이 클라우드로 나가면 저는 잘립니다' / '340개사를 한 화면에서' & OP-07·04 \\ \hline
\fi
\end{tabularx}}
\landscapeoff

%% ------------------------------------------------------------------ 4
\secbar{4. 역할군별 종합}
\guide{역할군마다 공통 패턴 2~3개와 서로 상충하는 것(검사원 vs 품질팀장, 공장장 vs 1차사)을 적는다.}
{\footnotesize
\begin{tabularx}{\linewidth}{|L{26mm}|Y|L{60mm}|}\hline
\hd{역할군} & \hd{공통 패턴} & \hd{상충·주의} \\ \hline
\ifwbblank
\rd{11mm}검사원 &  &  \\ \hline
\rd{11mm}품질팀장 &  &  \\ \hline
\rd{11mm}공장장 &  &  \\ \hline
\rd{11mm}대표 &  &  \\ \hline
\rd{11mm}1차사 SQE &  &  \\ \hline
\else
검사원(11) & 라인이 도는 동안은 못 적는다(9/11) · 이미 사진을 찍는다 · 도구 실패 경험 2건 & 시간 단축은 환영, 판정 대체는 침묵(I-02·11·15). 배석(I-18) 신뢰도 낮음 \\ \hline
품질팀장(8) & 통보 지연 원인은 자료 수집(4/8) · 8D는 붙이는 것(3/8) · 포털 규격 변경(2/8) & 팀장은 근거·속도, 검사원은 시간 — 같은 기능이 통제 vs 감시 \\ \hline
공장장(4) & 판정 일관성(2/4) · 데이터 활용(1/4) · IT 담당 부재(1/4) & '빨리보다 똑같이' — 예산 승인자이나 사용자 아님 \\ \hline
대표(3) & 클레임 근거(2/3) · 1차사 요구 시 도입(1/3) · '시스템은 사치'(1/3) & 지불 결정은 클레임·1차사 압력에서 — 내부 효율로는 안 움직임 \\ \hline
1차사 SQE(1) & 도면 클라우드 금지 · 340개사 한 화면 & 최대 블로커이자 최대 채널 — ㉕ 첫 결정 \\ \hline
\fi
\end{tabularx}}

%% ------------------------------------------------------------------ 5
\secbar{5. 기회 문장 목록 — 해법 명사 검증}
\guide{'검사원이 하루 31분을 이미 적힌 값을 다시 적는 데 쓴다'는 기회, '태블릿이 필요하다'는 무효. 검증 열에 ✓/✗.}
{\footnotesize
\begin{tabularx}{\linewidth}{|L{14mm}|Y|L{30mm}|L{22mm}|L{22mm}|}\hline
\hd{OP-ID} & \hd{기회 문장(상황 + 진전 + 장애)} & \hd{근거 인터뷰 ID} & \hd{빈도/강도} & \hd{해법 명사 검증} \\ \hline
\ifwbblank
\rd{9mm} &  &  &  &  \\ \hline
\rd{9mm} &  &  &  &  \\ \hline
\rd{9mm} &  &  &  &  \\ \hline
\rd{9mm} &  &  &  &  \\ \hline
\rd{9mm} &  &  &  &  \\ \hline
\rd{9mm} &  &  &  &  \\ \hline
\rd{9mm} &  &  &  &  \\ \hline
\rd{9mm} &  &  &  &  \\ \hline
\else
OP-01 & 검사원이 라인이 도는 동안 적은 값을 하루 31분 동안 다시 적는다(종이 → 엑셀 → 포털) & I-01·02·04·07·10·13·14·18·20·21·24·25 & 12(9/2/0/1/0) / 3 & ✓ \\ \hline
OP-02 & 부적합이 나와도 사진·근거·기록을 모으는 데 이틀이 걸려 1차사 통보가 평균 3.7일 늦는다 & I-03·08·19·24·25 & 5(1/4/0/0/0) / 3 & ✓ \\ \hline
OP-03 & 품질팀장이 8D 한 건에 11.5시간을 쓰고 그 대부분이 옮겨 적는 시간이다 & I-03·05·16 & 3(0/3/0/0/0) / 3 & ✓ \\ \hline
OP-04 & 클레임·감사 때 '누가 언제 무엇을 보고 판정했는가'를 찾을 수 없다 & I-09·11·12·23·27 & 5(1/1/0/2/1) / 2 & ✓ \\ \hline
OP-05 & 신입·교대 검사원의 판정이 숙련자와 다르고 그 차이를 아무도 재지 않는다 & I-02·06·11·15·26 & 5(3/0/2/0/0) / 2 & ✓ \\ \hline
OP-06 & 현장의 도면·검사기준 버전이 최신과 다르고 검사원은 그것을 알 수 없다 & I-04·15·20 & 3(3/0/0/0/0) / 2 & ✓ \\ \hline
OP-07 & 1차사 포털 규격이 바뀔 때마다 다시 배우고, 도면이 밖으로 나가면 SQ 등급이 강등된다 & I-08·22·23·27 & 4(0/1/1/1/1) / 2 & ✓ \\ \hline
OP-08 & 검사 데이터로 공정을 개선하고 싶지만 데이터가 없다 & I-14·17 & 2(0/0/1/1/0) / 1 & ✓ — 기각 후보(바람) \\ \hline
(무효) & 사진으로 기록해야 한다 / 태블릿이 필요하다 / AI 판정이 필요하다 & — & — & ✗ 해법 명사 → OP-01·05로 \\ \hline
\fi
\end{tabularx}}

%% ------------------------------------------------------------------ 6
\secbar{6. 관찰 vs 자기 보고의 괴리}
\guide{"다 해봤어요"라고 말한 사람이 실제로 어떻게 행동했는가. 관리자 배석 여부를 적는다. 이 표가 Day7에 로그와 대조된다.}
{\footnotesize
\begin{tabularx}{\linewidth}{|L{16mm}|Y|Y|L{50mm}|}\hline
\hd{인터뷰 ID} & \hd{자기 보고} & \hd{관찰·로그·제3자 진술} & \hd{해석·다음 질문} \\ \hline
\ifwbblank
\rd{12mm} &  &  &  \\ \hline
\rd{12mm} &  &  &  \\ \hline
\rd{12mm} &  &  &  \\ \hline
\rd{12mm} &  &  &  \\ \hline
\else
I-02 박선재 & '다 해봤어요, 결국 종이에' — 도구 무용론 & 배석 전 '느려서' / 배석 후 '괜찮았다' · 오후분 종이 메모 & 판정 기록이 남는 것에 대한 반응? → A-05. 배석 없이 라인에서 재인터뷰(X-01) \\ \hline
I-10 한빛금속 & '입력이 종이보다 느렸어요' & 태블릿이 서랍에, 배터리 없음 — 충전 담당 없음 & 속도인가 운영인가 → A-09 \\ \hline
I-18 대양정밀 & '괜찮아요, 할 만해요' & 팀장 배석, 12분 & 신뢰도 낮음 — 빈도 가중 0. 현장 재방문 \\ \hline
I-24 미래오토 & '사진은 찍어요, 그다음이 문제' & 검사 중 폰 촬영 → 카톡 → 팀장 & 행동을 바꿀 필요가 없는 세그먼트 — A-01 부분 증거 \\ \hline
\fi
\end{tabularx}}

%% ------------------------------------------------------------------ 7
\secbar{7. 다음 인터뷰 계획}
\guide{누구를 더 만나야 하는가(빠진 역할군·상충을 풀 사람·LOI 후보), 언제까지, 무엇을 확인하려고.}
\begin{fieldbox}
\ifwbblank
\lines{5}{9.5mm}
\else
1. 검사원 재인터뷰 3건 — 배석 없이 라인에서(I-02·11·18), X-01 기간에.
2. 1차사 SQE 2명 추가 — 도면 클라우드 금지가 H사만의 정책인지(A-07 분모).
3. 사출·프레스 검사원 3건 — 절삭·단조 편향 보정.
4. 대표 2건 — 클레임 경험 있는 회사만(A-03).
5. 30인 미만 공장 2건 — 초기 SAM 하한의 근거.
\fi
\end{fieldbox}

\end{document}